\documentclass[11pt]{article}
\usepackage[a4paper,margin=23mm]{geometry}
\usepackage[T1]{fontenc}
\usepackage{lmodern,amsmath,amssymb,amsthm,booktabs,longtable,array,microtype}
\usepackage{xcolor}
\usepackage[colorlinks=true,linkcolor=blue!50!black,urlcolor=blue!50!black]{hyperref}
\newtheorem{theorem}{Theorem}
\newtheorem{lemma}[theorem]{Lemma}
\newcommand{\R}[1]{\texttt{R12\_1\_#1}}
\newcommand{\FPdim}{\operatorname{FPdim}}
\newcommand{\Rep}{\operatorname{Rep}}
\newcommand{\C}{\mathbb C}
\newcommand{\Z}{\mathbb Z}
\newcommand{\Q}{\mathbb Q}
\setlength{\parskip}{4pt}
\setlength{\parindent}{0pt}
\title{Rank twelve, multiplicity one:\\complex categorification}
\author{Fusion-ring census research notes}
\date{26 October 2026}
\begin{document}
\maketitle
\begin{abstract}
We classify complex categorifiability of the 632 supplied rank-twelve,
multiplicity-one fusion rings: 480 are excluded and 152 admit models.
The models include 151 unitary realizations and one additional
complex/pseudo-unitary realization whose unitarity is not proved.
The exact ledger binds every decision to the original full multiplication
tensor, categorical proof and reproducible certificate. The complete
proof archive retains the earlier family reductions and constructions;
this manuscript gives a unified reading guide and the final theoretical
arguments. A new equality-case refinement of the dimension-block conductor
bound excludes the final two commutative rings. An algebraic near-group
quotient realizes the last noncommutative ring and yields a uniform family.
Enumeration completeness is producer-reported; exact tensor validity and
based-isomorphism uniqueness were independently checked computationally,
without repeating the full search. The proofs and computations are
AI-generated and await independent human verification and formalization.
\end{abstract}

\section*{Classification and proof guide}
The inputs are the 632 physical lines of the supplied rank-twelve
multiplicity-one tensor file, with permanent labels
\texttt{R12\_1\_1},\ldots,\texttt{R12\_1\_632}.
A model means a fusion category over $\mathbb C$ together with an exact
based-ring identification; it does not classify associators.
An exclusion rules out arbitrary complex fusion categories, with any
pivotal, spherical or pseudo-unitary reduction justified separately.

The complete accepted ledger, source tensors, code, proofs and separate
assistant audits are in the
\href{https://github.com/sebastienpalcoux/grothendieck-ring-classification/tree/main/evidence/rank12_m1_20261026}{rank-twelve proof archive}.
Its top-level ledger is authoritative. Historical ledgers retain their
original, explicitly dated scope; their former open counts are not current.
Every earlier settled decision is retained. The first stage supplies
393 exclusions and 127 unitary models. The three subsequent family
continuations add respectively $(32,14)$, $(13,7)$ and $(40,3)$
exclusions and unitary models. The final arguments below add two
exclusions and one complex/pseudo-unitary model. Thus
\[
393+32+13+40+2=480,\qquad 127+14+7+3+1=152.
\]
No status follows merely from failure to find a model, a numerical
optimization, or a conditional modular datum.

The archive's four earlier reports give the full proofs and accepted
replays for lower-rank subring inheritance, cyclic isotypic integrality,
formal-codegree and dimension-order obstructions, paired stabilizers,
pointed Morita extensions, cohomology and relative-algebra reconstructions.
The model certificates use standard quantum-group and group-theoretical
categories, pointed and TY extensions, coherent equivariantizations,
central zests and actual Morita duals. Each construction checks its
categorical hypotheses before identifying all fusion coefficients.
The accepted replay instructions and frozen source manifests accompany
each report. The current verifier binds the entire preceding packet
byte for byte; its optional baseline replay executes the full earlier
audit. The final three arguments are proved in detail below.



\section{Decisions and common structure}
\begin{center}
\begin{tabular}{lll}\toprule
Ring & Decision & Reason\\\midrule
\R{561}&Excluded over $\C$&Central $C_3$ and conductor saturation\\
\R{628}&Excluded over $\C$&The same argument\\
\R{623}&Modeled over $\C$&Cyclic order-$36$ near-group quotient\\\bottomrule
\end{tabular}
\end{center}
All three have six invertible objects forming $G=C_6=\langle g\rangle$
and six noninvertibles forming a free $G$-orbit. Put $z=g^3$ and
\[
 d=3+\sqrt{10},\qquad d^2=1+6d,\qquad
 D=6(1+d^2)=120+36\sqrt{10}.
\]
Their presentations are
\[
\begin{array}{c|cc}
 &\rho g&\rho^2\\\hline
\R{561}&g\rho&1+\sum_{h\in G}h\rho\\
\R{628}&g\rho&z+\sum_{h\in G}h\rho\\
\R{623}&g^{-1}\rho&z+\sum_{h\in G}h\rho.
\end{array}
\]
The universal grading of each ring is trivial: the products $XX^*$
generate the entire based ring. The complete tensors, rather than only
these abbreviated presentations, are bound in the accompanying packet.
The $629$ previously settled decisions are retained unchanged. Together with the three decisions below, they give the full classification of the supplied census.

\section{The central \texorpdfstring{$C_3$}{C3} premise for the two commutative rings}
We recall the categorical premise from the preceding continuation,
with its exact verification included again. A hypothetical complex
categorification can be replaced by a pseudo-unitary Galois conjugate
using the two-dimension theorem \cite[Lemma 3.1]{Sch23}.
Let $H=\langle g^2\rangle\cong C_3$. Conjugation by the actual invertible
objects gives a coherent $H$-action, even when the pointed associator
is nontrivial: its scalar occurs with its inverse on the dual factor.
Form the genuine equivariantization $E=\mathcal C^H$.

Every source simple has three equivariant lifts, since its stabilizer is
cyclic. The pointed group $A$ of $E$ is a central extension of $C_6$ by
$K=\widehat H$, hence is $C_{18}$ or $C_6\times C_3$. The noninvertible
action is free. Choosing a lift $r$ of $\rho$ gives
\[
 r^2=t'+\sum_{s\in S}sr,
\]
where $S$ selects one element in every fiber $A\to C_6$. Associativity
forces the lifted right action to be trivial; duality imposes
$S^{-1}=t'^{-1}S$. There are exactly $72$ possibilities across the two
rings and the two pointed groups. For each character $\lambda$ of $A$,
write $s_\lambda=\sum_{s\in S}\lambda(s)$ and
$\epsilon=\lambda(t')$. The two associated formal codegrees satisfy
\[
 f_+f_-=18^2(4+s_\lambda^2/\epsilon).
\]
Central induction requires $(3D)/f_\pm$ to be algebraic integers.
Exact cyclotomic norms eliminate $66$ sections: a denominator has a prime
factor outside $\{2,3,5\}$, while the numerator is a unit at that prime.
The six surviving sections are subgroup complements $B$ of $K$ in
$C_6\times C_3$, three for each ring. The replay enumerates this entire
small domain; it is not a sampled search.

For a surviving section, $E_{\mathrm{ad}}$ has simples $B\cup Br$.
The forgetful functor $E_{\mathrm{ad}}\to\mathcal C$ sends these
bijectively to distinct simple objects, and is therefore a tensor
equivalence. Equivariant structures for the inner action provide central
lifts of every $h\in H$; equivalently, use the reverse central embedding
from the center of the pointed $H$-subcategory into the center of its
relative center. Select the unique lift lying in $E_{\mathrm{ad}}$.
Half-braidings on a fixed invertible are unique because
$U(E_{\mathrm{ad}})=1$, so the selected lifts close under tensor product
as central objects. Transporting them gives a fully faithful braided
$C_3$ in $Z(\mathcal C)$ whose forgetful image is the original $H$.

This argument proves an actual central lift, not just commutativity of
the fusion ring. Its braiding is either Tannakian or one of the two
nondegenerate forms $q(a)=\zeta_3^{ra^2}$, $r=1,2$.
The next section excludes all three choices.

\section{One conductor argument excludes both rings}
\subsection{The equality refinement}
For a spherical fusion category let $N$ be its Frobenius--Schur exponent.
Its center has a canonical unitary modular representation
\[
 \varrho:\mathrm{SL}_2(\Z/N)\longrightarrow U(V),
 \qquad \varrho(s)=S,\quad\varrho(t)=T,
\]
where $S$ is normalized. This unitarity concerns the matrices of the
modular representation; it does not assume that the original tensor
category is unitarizable. The dimension-block bound \cite{Sch26}
associates to a dimension-field character $\chi$ a twist-fixed vector
$w_\chi$, with diagonal-torus character $\chi$, and bounds its
$S$-matrix coefficient by $\mathfrak f(\chi)^{-1/2}\|w_\chi\|^2$.

\begin{lemma}[Auxiliary-prime rigidity at equality]
Suppose this bound is attained. For every prime $p\mid N$ with
$p\nmid\mathfrak f(\chi)$, the vector $w_\chi$ is fixed by the entire
local factor $\mathrm{SL}_2(\Z/p^{v_p(N)})$.
\end{lemma}
\begin{proof}
Use the Chinese remainder decomposition of the modular representation.
For each prime $\ell\mid N$, let $U_\ell$ be the upper unipotent group,
$P_\ell$ its fixed-space projection, and
$B_\ell=P_\ell\varrho(s_\ell)P_\ell$.
The local estimate in \cite[Lemma 4.1]{Sch26} bounds the product of
the compressions at primes other than $p$ by
$c=\mathfrak f(\chi)^{-1/2}$ on the relevant character space;
$\|B_p\|\leq1$. The operators commute. Equality gives
\[
 c\|w\|^2
 =|\langle w_{\bar\chi},B_{\mathrm{rest}}B_pw\rangle|
 \leq c\|w\|\,\|B_pw\|\leq c\|w\|^2.
\]
Here $\|w_{\bar\chi}\|=\|w\|$.
Thus $\|P_p\varrho(s_p)w\|=\|w\|$, so $\varrho(s_p)w$ is
$U_p$-fixed. The vector $w$ is fixed by $U_p$ and by its opposite
$s_p^{-1}U_ps_p$. These elementary unipotent subgroups generate
$\mathrm{SL}_2(\Z/p^{v_p(N)})$. This works on direct sums of tensor
product constituents as well; $w$ need not be a pure tensor in the
Chinese remainder decomposition.
\end{proof}

In our two rings the simple-dimension and global-dimension fields are
both $\Q(\sqrt{10})$. Its nontrivial character has conductor $40$.
Writing $D'=120-36\sqrt{10}$ gives
\[
 DD'=1440,\qquad D/D'=d^2,\qquad
 \left(\frac1D-\frac1{D'}\right)^2=\frac1{40}.
\]
Let $A_\sigma$ be the simple summand of $I(1)$ associated to the
conjugate dimension character. Then $\dim A_\sigma=d^2$,
$\theta_{A_\sigma}=1$, and
\[
 w=e_1-e_{A_\sigma},\qquad
 |\langle w,Sw\rangle|=\frac{\|w\|^2}{\sqrt{40}}.
\]
The lemma therefore applies at the prime $3$ whenever $3\mid N$.

\subsection{Tannakian and metric alternatives}
If the central $C_3$ is Tannakian, the actual regular algebra gives
tensor de-equivariantization of dimension $D/3$. Its four simple
dimensions are $1,1,d,d$, because the $C_3$-action on source simples is
free. A pseudo-unitary Galois conjugate of this quotient would violate
the same dimension-block bound:
\[
 \left(\frac3D-\frac3{D'}\right)^2
 =\frac9{40}>\frac1{40}.
\]
No rank-four classification is needed.

Otherwise the central $C_3$ is a nondegenerate pointed category
$\mathcal B$. Apply modular decomposition \emph{inside the braided
center}:
\[
 Z(\mathcal C)\simeq\mathcal B\boxtimes\mathcal D.
\]
The induction-character identity gives
$S^{\mathrm{un}}_{A_\sigma,h}=\dim(A_\sigma)$ for $h\in C_3$,
since both dimension characters send its underlying invertible to $1$.
As $h$ is invertible, its monodromy with $A_\sigma$ is therefore $1$.
Thus $A_\sigma\in\mathcal D$ and
\[
 w=e_0\otimes(e_1-e_Y)
\]
is a nonzero pure tensor for this modular decomposition.

The pointed twists force $3\mid N$. Put $n=N/3^{v_3(N)}$, so
$3\nmid n$. The element $L=st^ns^{-1}$ lies entirely in the local
$3$-factor, because $t^n$ is the identity at every other prime power.
The equality lemma forces $\varrho(L)w=w$.
Normalized modular matrices factor exactly, so
\[
 \varrho(L)=(S_{\mathcal B}T_{\mathcal B}^nS_{\mathcal B}^{-1})
 \otimes(S_{\mathcal D}T_{\mathcal D}^nS_{\mathcal D}^{-1}).
\]
For $q(a)=\zeta_3^{ra^2}$, set $u=\zeta_3^{rn}\ne1$.
The Fourier matrix of $\mathcal B$ gives
\[
 S_{\mathcal B}T_{\mathcal B}^nS_{\mathcal B}^{-1}e_0
 =\frac13(1+2u,\,1-u,\,1-u).
\]
This is not proportional to $e_0$. The second tensor factor is
invertible and sends $e_1-e_Y$ to a nonzero vector. Hence
$\varrho(L)w\ne w$, a contradiction.

\begin{theorem}
Neither \R{561} nor \R{628} is the Grothendieck ring of a complex
fusion category.
\end{theorem}
The proof uses no bound on the unknown exponent $N$, no modular
representation phase choice for either factor, no ambient braiding on
$\mathcal C$, and no categorical unitarity assumption. The additional
equality argument is developed here from the proof of \cite{Sch26}.

\section{A complex model for \R{623}}
\subsection{An existing algebraic precursor}
Let $G_0=C_{36}$, additively. The algebraic construction
\cite[Theorems 2.3 and 2.10]{TPE} supplies an actual spherical near-group
category $\mathcal N$ with simple objects $G_0\cup\{X\}$ and
\[
 gX=Xg=X,\qquad X^2=\sum_{g\in G_0}g+36X,
 \qquad\FPdim(X)=18+6\sqrt{10}.
\]
The construction first has negative categorical dimension for $X$;
its positive spherical Galois conjugate is also available. Its quadratic
data are
\[
 \chi(g,k)=e^{2\pi i gk/36},\qquad
 a(g)=e^{\pi i g(36-g)/36}.
\]
In the conventions of \cite{Izumi}, the parameters needed below are
$b=\chi^{-1}$ and $a_I=a^{-1}$, by \cite[Lemma 2.14]{TPE}.
We use the spherical existence theorem, not the conjecture about
unitarity at order $36$.

\subsection{The actual central algebra}
Take $H=6\Z/36\Z$, of order $6$. Then
\[
 b|_{H\times H}=1,\qquad H^\perp=H,
 \qquad a_I(6j)=(-1)^j.
\]
The central-invertible calculation of \cite[Section 4.3]{EG}, also
written explicitly in the proof of \cite[Theorem 13.4]{Izumi}, applies
algebraically to the explicit mixed associators of this precursor.
For clarity, this is not an assertion that every spherical near-group
category belongs to a unitary classification. In the reconstructed
mixed-vertex basis $v_t$, the three group insertions have the form
\[
 Z(h)v_t=\chi(h,t)v_t,\quad X(h)v_t=v_{t+h},\quad
 Y(h)v_t=\frac{a(t)}{a(t-h)}v_{t-h}.
\]
Their product satisfies
\[
 X(h)Y(h)Z(h)=a(h)^{-1}I.
\]
The other identities needed for the invertible half-braiding are
bicharacter multiplication and $a(h)^2\chi(h,h)=1$.
For $h\in H$, the candidate mixed crossing uses $a(h)$;
$a(h)^2=1$ checks the duality channel and the displayed product checks
the trivalent channel. All maps are inverse mixed vertices or bilinear
dual maps in the algebraic reconstruction \cite[Section 4.7]{TPE}.
There is no positivity condition and no all-$X$ associator block in
these hexagons. The detailed convention audit is retained in the packet.

The central lifts restrict to a Tannakian $H$: their mutual braiding
is $b|_{H\times H}=1$. They close coherently because $U(\mathcal N)=1$
makes half-braidings on invertibles unique. The pointed associator is
trivial, since the fixed simple $X$ is a rank-one module category for
the pointed subcategory. Consequently
\[
 A=\bigoplus_{h\in H}h
\]
with regular group multiplication is an actual connected commutative
separable algebra in $Z(\mathcal N)$. Its module category
$\mathcal M=\mathcal N_A$ is a fusion category.

\subsection{Fusion rules without a unitary reconstruction assumption}
The quotient has six invertibles $u_j$ indexed by $G_0/H=C_6$.
The induced $X$ splits into six distinct simple summands using the
primitive Fourier idempotents of $\C[H]$. The perfect pairing
$G_0/H\to\widehat H$ makes their pointed action free. Left and right
multiplication shift the character with opposite signs, giving
\[
 \sigma u_j=u_{-j}\sigma.
\]
The pointed-channel calculation in the isotropic orbifold
\cite[Theorem 12.5]{Izumi} is the algebraic character average
\[
 \frac16\sum_{h\in H}a_I(h)b(h,g)
 =\begin{cases}1&g\equiv3\pmod6,\\0&\text{otherwise.}\end{cases}
\]
This uses only the mixed vertices and their inverse dual maps.
Thus the unique pointed term of $\sigma^2$ is $z=u_3$, and
$\sigma^*=z\sigma$.

The remaining coefficients require no transported operator-algebra
calculation. Write
\[
 \sigma^2=z+\sum_{j=0}^5m_j u_j\sigma.
\]
Associativity with the invertibles forces $m_{j+2}=m_j$.
The noninvertible dimension is $d$, so $\sum_jm_j=6$.
The inversion rule gives $(u_j\sigma)^*=u_{j+3}\sigma$,
while $(\sigma^2)^*=\sigma^2$. Duality therefore forces
$m_{j+3}=m_j$. Shifts by $2$ and $3$ generate $C_6$, so all six
coefficients equal $1$. We obtain exactly
\[
 \sigma g=g^{-1}\sigma,
 \qquad\sigma^2=z+\sum_{g\in C_6}g\sigma.
\]
The coordinate-to-census bijection, listing
$1,g,\ldots,g^5,\sigma,g\sigma,\ldots,g^5\sigma$, is
\[
 [0,1,10,11,9,2,3,8,5,4,7,6].
\]
Two independent implementations check all $1728$ structure constants,
duality, and the exact quadratic-form specialization.

\begin{theorem}
The ring \R{623} has a complex categorification and a pseudo-unitary
categorification. The present argument does not establish unitarity.
\end{theorem}
The strict isotropic statement above is essential: a version requiring
$a_I|_H=1$ would miss the shift $z$ and give the wrong ring.

\subsection{A uniform family}
The same construction works for every $k\equiv2\pmod4$. Start from
the cyclic near-group category $C_{k^2}+k^2$ and take
$H=k\Z/k^2\Z$. The pairing is trivial on $H$ and $H=H^\perp$,
while $a_I(kj)=(-1)^j$. The quotient therefore has $k$ invertibles,
$k$ noninvertibles, inversion mixed action, and pointed square term
$g^{k/2}$. Associativity makes the noninvertible coefficients constant
on parity classes; duality shifts the index by the odd integer $k/2$
and equates those classes. Their sum is $k$, hence every coefficient
is $1$. Thus the construction supplies, uniformly over $\C$,
\[
 \rho g=g^{-1}\rho,\qquad
 \rho^2=g^{k/2}+\sum_{j=0}^{k-1}g^j\rho,
 \qquad k\equiv2\pmod4.
\]
The case $k=6$ is \R{623}. This corollary inherits the same
distinction between complex existence and unproved unitarity.

\section{Evidence, totals, and limits}
The accompanying archive contains the updated $632$-ring ledger, the
three complete input tensors, the new proofs, separate proof reviews,
and portable exact checks. It preserves the entire preceding proof
packet. The current checks replay the forced-central-$C_3$ premise,
the rational and cyclotomic equality calculations, the algebraic
quotient data, and both independent full-tensor identifications.
There is no numerical associator search or long classification run.

\begin{center}
\begin{tabular}{lrrr}\toprule
Stage&Excluded&Modeled&Unresolved\\\midrule
Previous ledger&478&151&3\\
This continuation&2&1&$-3$\\
Final ledger&\textbf{480}&\textbf{152}&\textbf{0}\\\bottomrule
\end{tabular}
\end{center}
The new source inputs are the pinned preprints \cite{Sch26,TPE}.
The exact checks certify their application and the census bindings;
they are not replacements for the categorical proofs or an independent
reconstruction of the analytic precursor theorem. The model is
``new'' only as an explicit realization added to this project, with no
claim of novelty up to tensor equivalence or exoticness.
The companion census repository records enumeration separately from these categorification decisions.

\begin{thebibliography}{9}
\bibitem{Sch23} A. Schopieray, \emph{Categorification of integral group
rings extended by one dimension}, J. London Math. Soc. 108 (2023),
1617--1641. \href{https://doi.org/10.1112/jlms.12789}{doi:10.1112/jlms.12789}.
\bibitem{Sch26} A. Schopieray, \emph{The dimension block of a spherical
fusion category}, \href{https://arxiv.org/abs/2609.35287v1}{arXiv:2609.35287v1},
28 September 2026. Lemmas 3.1, 4.1 and Theorem 5.1.
\bibitem{TPE} H. He, Z. Liu, F. Lu, S. Palcoux and Y. Ren,
\emph{Triangular prism equations for near-group categories: algebraic
solutions and a unitary conjecture},
\href{https://arxiv.org/abs/2610.11652v1}{arXiv:2610.11652v1},
8 October 2026. Theorems 2.3, 2.10, 5.4; Lemmas 2.14, 4.14;
Corollary 6.6.
\bibitem{EG} D. E. Evans and T. Gannon, \emph{Near-group fusion
categories and their doubles},
\href{https://arxiv.org/abs/1208.1500}{arXiv:1208.1500},
Section 4.3, equations (4.48)--(4.50).
\bibitem{Izumi} M. Izumi, \emph{A Cuntz algebra approach to the
classification of near-group categories},
\href{https://arxiv.org/abs/1512.04288v2}{arXiv:1512.04288v2},
Section 12.1, Theorem 12.5, and the proof of Theorem 13.4.
\end{thebibliography}
\end{document}
